\subsection{Kritische Würdigung}
\label{subsec:kritische-wuerdigung}

Die technische Fallstudie untersucht das konkrete Template detailliert, ist
aber nur begrenzt generalisierbar. Projektunterlagen belegen die Implementierung,
nicht deren unabhängige Evaluation; Abnahme und lokale Prüfung stammen
ebenfalls aus dem Projektkontext. Die Verbindung von Entwicklerperspektive,
Selbstdokumentation und Bewertung kann die Evidenzauswahl beeinflussen, ohne
commitgebundene Befunde aufzuheben. Wissenschaftliche Literatur ordnet
allgemeine Konzepte ein, offizielle Dokumentation die Technologien; eine
unabhängige Replikation ersetzen beide nicht.

Die Testevidenz ist heterogen: Einer früheren, überwiegend Linux-basierten
Abnahme stehen gemischte aktuelle Remote-Ergebnisse gegenüber. Jeder Test und
Build belegt nur sein Kriterium; ohne Coverage-Messung, Compose-Start und
durchgehend erfolgreiche Plattformläufe entsteht kein umfassender Qualitäts-
oder Produktionsnachweis. Unsignierte Pakete ersetzen zudem keine nativen
Produktreleases einschließlich Signierung und Notarisierung.

Übertragbar erscheinen die Prinzipien zentraler Baseline, deklarativer
Variabilität und gemeinsamer Tooling-Einstiege; projektspezifisch bleiben
Technologien, Profile und commitgebundene Evidenz. Produktivitätspotenzial und
Projekteignung wurden qualitativ hergeleitet. Zeitmessungen, Vergleichsgruppen,
reale Kundenfallprüfungen und eine wirtschaftliche Bewertung fehlen.
